\documentclass[a4paper,12pt,titlepage,finall]{article}

\usepackage[T1,T2A]{fontenc}     % форматы шрифтов
\usepackage[utf8x]{inputenc}     % кодировка символов, используемая в данном файле
\usepackage[russian]{babel}      % пакет русификации
\usepackage{tikz}                % для создания иллюстраций
\usepackage{pgfplots}            % для вывода графиков функций
\usepackage{geometry}		 % для настройки размера полей
\usepackage{indentfirst}         % для отступа в первом абзаце секции
\usepackage{multirow}            % для таблицы с результатами

% выбираем размер листа А4, все поля ставим по 3см
\geometry{a4paper,left=30mm,top=30mm,bottom=30mm,right=30mm}

\setcounter{secnumdepth}{0}      % отключаем нумерацию секций

\usepgfplotslibrary{fillbetween} % для изображения областей на графиках

\begin{document}
% Титульный лист
\begin{titlepage}
    \begin{center}
	{\small \sc Московский государственный университет \\имени М.~В.~Ломоносова\\
	Факультет вычислительной математики и кибернетики\\}
	\vfill
	{\Large \sc Отчет по заданию №1}\\
	~\\
	{\large \bf <<Методы сортировки>>}\\ 
	~\\
	{\large \bf Вариант 1 / 2 / 3 / 4}
    \end{center}
    \begin{flushright}
	\vfill {Выполнил:\\
	студент 104 группы\\
	Денисов~М.~А.\\
	~\\
	Преподаватель:\\
	Гуляев~Д.~А.}
    \end{flushright}
    \begin{center}
	\vfill
	{\small Москва\\2026}
    \end{center}
\end{titlepage}

% Автоматически генерируем оглавление на отдельной странице
\tableofcontents
\newpage

\section{Постановка задачи}

В рамках данной работы решается задача экспериментального исследования эффективности двух алгоритмов сортировки массивов. Согласно варианту, исследованию подлежат метод сортировки выбором - Selection sort - и метод пирамидальной сортировки – Heapsort. Оба метода предназначены для сортировки элементов одномерного массива по возрастанию.

Исследование заключается в оценке вычислительной сложности алгоритмов путем подсчета двух операций:
\begin{itemize}
\item Число сравнений элементов массива
\item Число перемещений (обменов) элементов в процессе сортировки.
\end{itemize}
Для объективности эксперимента тестовые массивы будут генерироваться с разной исходной упорядоченностью. Генерация массивов выполняется отдельной функцией, которая в зависимости от параметра создает следующие конфигурации:
\begin{itemize}
\item Массив, упорядоченный по возрастанию
\item Массив, упорядоченный по убыванию
\item Массив со случайной расстановкой элементов
\item Массив с иной случайной расстановкой элементов. 
\end{itemize}

Сравнение методов сортировки производится на одних и тех же исходных массивах. Эксперименты проводятся для массивов различной длины. Память будет динамически выделяться под размеры n = 10, 100, 1000, 10000.

Итоговые данные усредняются и оформляются в виде таблиц, отражающих зависимость числа сравнений и перемещений от типа исходных данных и размера массива для каждого из исследуемых алгоритмов.

\newpage

\section{Результаты экспериментов}

\begin{table}[h]
\centering
\begin{tabular}{|c|c|c|c|c|c|c|c|}
    \hline
    \multirow{2}{*}{\textbf{n}} & \multirow{2}{*}{\textbf{Параметр}} & \multicolumn{4}{|c|}{\textbf{Номер сгенерированного массива}} & \textbf{Среднее} \\
    \cline{3-6}
    & & \parbox{1.5cm}{\centering 1} & \parbox{1.5cm}{\centering 2} & \parbox{1.5cm}{\centering 3} & \parbox{1.5cm}{\centering 4} & \textbf{значение} \\
    \hline
    \multirow{2}{*}{10} & Сравнения & 45 & 45 & 45 & 45 & 45.00 \\
    \cline{2-7}
                        & Перемещения & 0 & 5 & 6 & 7 & 4.50 \\
    \hline
    \multirow{2}{*}{100} & Сравнения & 4950 & 4950 & 4950 & 4950 & 4950.00 \\
    \cline{2-7}
                         & Перемещения & 0 & 50 & 94 & 96 & 60.00 \\
    \hline
    \multirow{2}{*}{1000} & Сравнения & 499500 & 499500 & 499500 & 499500 & 499500.00 \\
    \cline{2-7}
                          & Перемещения & 0 & 516 & 991 & 992 & 624.75 \\
    \hline
    \multirow{2}{*}{10000} & Сравнения & 49995000 & 49995000 & 49995000 & 49995000 & 49995000.00 \\
    \cline{2-7}
                           & Перемещения & 0 & 6078 & 9991 & 9996 & 6516.25 \\
    \hline
\end{tabular}
\caption{Результаты работы сортировки выбором (Selection Sort)}
\end{table}

\begin{table}[h]
\centering
\begin{tabular}{|c|c|c|c|c|c|c|c|}
    \hline
    \multirow{2}{*}{\textbf{n}} & \multirow{2}{*}{\textbf{Параметр}} & \multicolumn{4}{|c|}{\textbf{Номер сгенерированного массива}} & \textbf{Среднее} \\
    \cline{3-6}
    & & \parbox{1.5cm}{\centering 1} & \parbox{1.5cm}{\centering 2} & \parbox{1.5cm}{\centering 3} & \parbox{1.5cm}{\centering 4} & \textbf{значение} \\
    \hline
    \multirow{2}{*}{10} & Сравнения & 41 & 35 & 41 & 41 & 39.50 \\
    \cline{2-7}
                        & Перемещения & 31 & 22 & 31 & 30 & 28.50 \\
    \hline
    \multirow{2}{*}{100} & Сравнения & 1081 & 944 & 1035 & 1021 & 1020.25 \\
    \cline{2-7}
                         & Перемещения & 641 & 517 & 588 & 578 & 581.00 \\
    \hline
    \multirow{2}{*}{1000} & Сравнения & 17580 & 15973 & 16883 & 16805 & 16810.25 \\
    \cline{2-7}
                          & Перемещения & 9708 & 8325 & 9124 & 9045 & 9050.50 \\
    \hline
    \multirow{2}{*}{10000} & Сравнения & 244428 & 226626 & 235369 & 235304 & 235431.75 \\
    \cline{2-7}
                           & Перемещения & 131236 & 116644 & 124203 & 124195 & 124069.50 \\
    \hline
\end{tabular}
\caption{Результаты работы пирамидальной сортировки (Heapsort)}
\end{table}

Сортировка выбором (Selection Sort):
\begin{itemize}
\item Сравнения строго постоянны ($n(n-1)/2$) для любых исходных данных — алгоритм всегда ищет минимум полным перебором.
\item Перемещения зависят от порядка: минимум (0) на отсортированном массиве, максимум ($\approx n$) на случайном, среднее ($n/2$) на обратном.
\item Подтверждена квадратичная сложность $O(n^2)$ по сравнениям и линейная $O(n)$ по обменам.
\end{itemize}

Пирамидальная сортировка (Heapsort):
\begin{itemize}
\item Сравнения и перемещения имеют сложность $O(n \log n)$ и варьируются в зависимости от данных.
\item Минимум операций — на обратном порядке, максимум — на отсортированном массиве.
\item Случайные данные дают промежуточные значения.
\end{itemize}

Сравнение при $n=10000$:
\begin{itemize}
\item Heapsort выполняет в 212 раз меньше сравнений, чем Selection Sort ($\approx 235\,000$ против $50$ млн).
\item По перемещениям Selection Sort эффективнее ($6\,500$ против $124\,000$), но это плата за логарифмическую сложность Heapsort.
\end{itemize}


\newpage

\section{Структура программы и спецификация функций}

Программа реализована на языке Си и состоит из набора функций, реализующих алгоритмы сортировки, генерацию тестовых массивов, подсчёт метрик и вывод результатов.

\subsection{Глобальные переменные}

В программе используются следующие глобальные переменные для подсчёта метрик:

\texttt{comp} — счётчик типа \texttt{int}, хранящий число сравнений элементов, выполненных в процессе сортировки.

\texttt{swap} — счётчик типа \texttt{int}, хранящий число обменов (перемещений) элементов.

\texttt{av\_comp\_sel}, \texttt{av\_comp\_heap} — накопительные счётчики для вычисления среднего числа сравнений для сортировки выбором и пирамидальной сортировки соответственно.

\texttt{av\_swap\_sel}, \texttt{av\_swap\_heap} — накопительные счётчики для вычисления среднего числа обменов.

\subsection{Функции сравнения}

\texttt{int compare(const void *a, const void *b)} — функция сравнения двух целых чисел. Возвращает отрицательное, нулевое или положительное значение в зависимости от того, меньше ли первый элемент второго. Используется для сортировки по возрастанию.

\texttt{int compare\_rev(const void *a, const void *b)} — функция сравнения для сортировки по убыванию. Возвращает значение, противоположное \texttt{compare}.

\subsection{Вспомогательные функции}

\texttt{void nullify()} — обнуляет глобальные счётчики \texttt{comp} и \texttt{swap} перед началом работы очередного алгоритма сортировки.

\texttt{void sw(int *a, int *b)} — выполняет обмен двух целочисленных элементов, переданных по указателю. При каждом обмене увеличивает счётчики \texttt{swap} и \texttt{av\_swap\_heap}.

\texttt{void arr\_pr(int *arr, int n)} — выводит на экран содержимое массива длины \texttt{n}. Используется для отладки при малых размерах массивов.

\subsection{Функции генерации массивов}

\texttt{void generate(int n, int type, int *ar1, int *ar2)} — универсальная функция, формирующая два идентичных массива \texttt{ar1} и \texttt{ar2} размера \texttt{n} в зависимости от параметра \texttt{type}:
\begin{itemize}
    \item \texttt{type = 1} — генерируется массив со случайными элементами, затем сортируется по возрастанию.
    \item \texttt{type = 2} — генерируется массив со случайными элементами, затем сортируется по убыванию.
    \item \texttt{type = 3} — генерируется массив со случайными элементами.
\end{itemize}
Функция также выводит заголовок с описанием типа массива и вызывает обе сортировки, выводя их результаты.

\subsection{Функции сортировки}

\texttt{void sel\_sort(int *arr, int n)} — реализует сортировку методом простого выбора (Selection Sort). Алгоритм на каждом шаге находит минимальный элемент в неотсортированной части массива и меняет его местами с первым элементом этой части. В процессе работы инкрементируются счётчики \texttt{comp} и \texttt{av\_comp\_sel} (при каждом сравнении), а также \texttt{swap} и \texttt{av\_swap\_sel} (при каждом обмене).

\texttt{void heapify(int *arr, int n, int i)} — вспомогательная рекурсивная функция для пирамидальной сортировки, обеспечивающая свойство невозрастающей пирамиды (max-heap) для поддерева с корнем \texttt{i}. Сравнивает корневой элемент с его левым и правым потомками, при необходимости выполняет обмен и рекурсивно вызывает себя для восстановления свойства пирамиды. Увеличивает счётчики \texttt{comp} и \texttt{av\_comp\_heap} при каждом сравнении.

\texttt{void heap\_sort(int *arr, int n)} — реализует пирамидальную сортировку (Heapsort). Сначала строит пирамиду из исходного массива с помощью функции \texttt{heapify}, затем последовательно извлекает максимальный элемент (корень), перемещая его в конец массива и восстанавливая пирамиду для оставшейся части. Использует функцию \texttt{sw} для обмена элементов.

\section{Отладка программы, тестирование функций}
Тестирование сортировок: при малых $N$ ($\leq 100$) вывод исходных и отсортированных массивов, сравнение с ожидаемым порядком, а также вывод массивов после каждой итерации сортировки.

\section{Анализ допущенных ошибок}
Ошибочный подсчет числа сравнений - в теле условия, а не вне.

\begin{raggedright}
\addcontentsline{toc}{section}{Список цитируемой литературы}
\begin{thebibliography}{99}
\bibitem{cs} Кормен Т., Лейзерсон Ч., Ривест Р, Штайн К. Алгоритмы: построение и анализ.
    Второе издание.~--- М.:<<Вильямс>>, 2005.
\end{thebibliography}
\end{raggedright}

\end{document}